% Part: methods
% Chapter: proofs
% Section: using-definitions

\documentclass[../../../include/open-logic-section]{subfiles}

\begin{document}

\olfileid[hi]{mod}{prf}{def}

\olsection{परिभाषाओं का उपयोग}

हमने कहा था कि प्रमाण में प्रयुक्त हो सकने वाली सभी परिभाषाओं से आपको
परिचित होना चाहिए और उन्हें ठीक प्रकार लागू करना आना चाहिए। यह वास्तव में
महत्त्वपूर्ण बात है और इसे कुछ विस्तार से देखना उपयोगी है। गुणों और संबंधों
को संक्षिप्त करने के लिए परिभाषाओं का उपयोग होता है, ताकि हम उनके बारे में
अधिक संक्षेप में बात कर सकें। प्रस्तुत संक्षिप्त रूप को \emph{परिभाष्य}
कहते हैं और जिसे वह संक्षिप्त करता है उसे \emph{परिभाषक}। प्रमाणों में हमें
प्रायः इस ओर लौटना पड़ता है कि परिभाष्य कैसे प्रस्तुत किया गया, क्योंकि
प्रमाण पूरा करने के लिए परिभाषक (वह लंबा रूप जिसका परिभाषित पद संक्षिप्त रूप
है) की तार्किक संरचना का उपयोग करना होता है। परिभाषाओं को खोलकर आप सुनिश्चित
करते हैं कि तार्किक क्रिया के मूल तक पहुँच रहे हैं।

एक उदाहरण से आरंभ करें। मान लें आप निम्न सिद्ध करना चाहते हैं:

\begin{prop}
किन्हीं समुच्चयों $A$ और $B$ के लिए $A\cup B=B\cup A$।
\end{prop}

प्रमाण आरंभ करने के लिए भी हमें जानना होगा कि दो समुच्चयों का समान होना क्या
है; अर्थात् हमें जानना होगा कि उस समीकरण में ``$=$'' का समुच्चयों के लिए
क्या अर्थ है। समुच्चय तब समान परिभाषित होते हैं जब उनके !!{element} समान
हों। इसलिए हमें यह परिभाषा खोलनी है:

\begin{defn}
समुच्चय $A$ और $B$ \emph{समान}, $A=B$, हैं तभी जब $A$ का प्रत्येक
!!{element}, $B$ का !!a{element} हो और इसके विपरीत भी।
\end{defn}

यह परिभाषा $A$ और $B$ को मनमाने समुच्चयों के स्थानधारकों के रूप में उपयोग
करती है। वह जिस चीज़ को परिभाषित करती है---\emph{परिभाष्य}---वह ``$A=B$''
अभिव्यक्ति है, और इसके लिए वह वह शर्त देती है जिसके अंतर्गत $A=B$ सत्य है।
यह शर्त---``$A$ का प्रत्येक !!{element}, $B$ का !!a{element} है और इसके
विपरीत भी''---\emph{परिभाषक} है।\footnote{इस विशेष स्थिति में---और बहुत
भ्रमकारी रूप से!---जब $A=B$, तो समुच्चय $A$ और $B$ एक ही समुच्चय हैं, यद्यपि
बाएँ और दाएँ उसके लिए भिन्न अक्षर उपयोग करते हैं। किंतु उस समुच्चय को निर्दिष्ट
करने के ढंग भिन्न हो सकते हैं और यही परिभाषा को अतुच्छ बनाता है।} परिभाषा
निर्धारित करती है कि $A=B$ ठीक तभी सत्य है जब शर्त सत्य हो।

परिभाषा लागू करते समय परिभाषा के $A$ और $B$ का मिलान संबंधित स्थिति से
करना होगा। हमारी स्थिति में इसका अर्थ है कि $A\cup B=B\cup A$ सत्य होने के
लिए प्रत्येक $z\in A\cup B$, $B\cup A$ में भी होना चाहिए और इसके विपरीत।
प्रतिज्ञप्ति में $A\cup B$ परिभाषा के $A$ की भूमिका निभाता है और $B\cup A$,
$B$ की। चूँकि $A$ और $B$ का उपयोग परिभाषा और सिद्ध की जा रही प्रतिज्ञप्ति
के कथन दोनों में, किंतु भिन्न उपयोगों में, होता है, सावधान रहना होगा कि दोनों
को न मिलाएँ। उदाहरणतः यह सोचना भूल होगी कि $A$ का प्रत्येक !!{element},
$B$ का !!a{element} है और इसके विपरीत दिखाकर प्रतिज्ञप्ति सिद्ध कर सकते
हैं---वह $A=B$ दिखाएगा, $A\cup B=B\cup A$ नहीं। (और चूँकि $A$ तथा $B$
कोई भी दो समुच्चय हो सकते हैं, आप अधिक आगे नहीं बढ़ेंगे, क्योंकि उनके बारे
में कुछ न माना हो तो वे भिन्न समुच्चय हो सकते हैं।)

प्रमाण में हम सम्मिलन जैसी समुच्चय-सैद्धांतिक धारणाओं से काम कर रहे हैं,
अतः प्रमाण कैसे आगे बढ़े यह समझने के लिए $\cup$ चिह्न का अर्थ भी जानना
होगा। और कभी-कभी किसी परिभाषा को खोलने से आगे खोलने योग्य परिभाषाएँ मिलती
हैं। उदाहरणतः $A\cup B$ को
$\Setabs{z}{z\in A\text{ अथवा }z\in B}$ के रूप में परिभाषित किया जाता है।
अतः यदि $x\in A\cup B$ सिद्ध करना हो, तो $\cup$ की परिभाषा खोलने से पता
चलता है कि $x\in\Setabs{z}{z\in A\text{ अथवा }z\in B}$ सिद्ध करना है। अब
यह भी याद रखना होगा कि $x\in\Setabs{z}{\dots z\dots}$ तभी जब
$\dots x\dots$। इसलिए $\Setabs{z}{\dots z\dots}$ संकेत की परिभाषा को और
खोलने पर दिखाना है: $x\in A$ अथवा $x\in B$। अतः ``$A\cup B$ का प्रत्येक
!!{element}, $B\cup A$ का भी !!a{element} है'' का वास्तविक अर्थ है:
``प्रत्येक $x$ के लिए यदि $x\in A$ अथवा $x\in B$, तो $x\in B$ अथवा
$x\in A$।'' प्रतिज्ञप्ति की परिभाषाओं को पूर्णतः खोलने पर दिखाना यह है:

\begin{prop}
किन्हीं समुच्चयों $A$ और $B$ के लिए: (a) प्रत्येक $x$ के लिए यदि $x\in A$
अथवा $x\in B$, तो $x\in B$ अथवा $x\in A$, और (b) प्रत्येक $x$ के लिए यदि
$x\in B$ अथवा $x\in A$, तो $x\in A$ अथवा $x\in B$।
\end{prop}

महत्त्वपूर्ण बात यह है कि परिभाषाओं को खोलना प्रमाण के निर्माण का आवश्यक
भाग है। इसे ठीक करना कभी-कभी कठिन होता है: परिभाषा के चरों और सिद्ध किए जा
रहे दावे के पदों में अंतर और मिलान सावधानी से करना होगा। सफल होने के लिए
आपको जानना होगा कि प्रश्न क्या पूछ रहा है और प्रश्न में प्रयुक्त सभी पदों का
अर्थ क्या है---प्रायः एक से अधिक परिभाषाएँ खोलनी पड़ेंगी। नीचे जैसे सरल
प्रमाणों में समाधान लगभग सीधे परिभाषाओं से ही निष्पन्न होता है। निस्संदेह
यह सदैव इतना सरल नहीं होगा।

\begin{prob}
मान लें आपसे $A\cap B\neq\emptyset$ सिद्ध करने को कहा गया है। यहाँ आने
वाली सभी परिभाषाएँ खोलिए, अर्थात् इसे ऐसे रूप में फिर कहिए जिसमें
``$\cap$'', ``$=$'' या ``$\emptyset$'' का उल्लेख न हो।
\end{prob}

\end{document}
